
%% -------------------------------------------------------------------
\subsection{\syncr{연구의 한계 (소규모 패널 \& 단일 $\kappa$)}}
\label{subsec:ch5_limitations}
%% -------------------------------------------------------------------

본 연구의 한계는 다음과 같다.

\fdel{\noindent\textbf{시뮬레이션 기반 검증의 한계.}
본 연구의 결과는 $n = 10{,}000$ 몬테카를로 시뮬레이션에 기반하며,
실제 장애 사건 데이터에 의한 검증은 수행되지 않았다.
시뮬레이션에서 사용한 입력 분포($k_O, O \sim \mathrm{Uniform}(0,1)$,
$U \sim N(3,1)$ 또는 $6 \cdot \mathrm{Pareto}(3)$)는 이론적
분석에 적합하나, 실제 조직의 관측성과 불확실성 분포를
완전히 반영하지 못할 수 있다.
향후 산업별 실제 장애 이력 데이터를 활용한 경험적 검증이 필요하다.}

\chg{전문가 패널의 규모 측면에서,}
\chg{정량적과 정성적} 전문가 검증을 5명의 소규모 패널로 수행하였다.
S-CVI/Ave $= 0.985$, Cronbach's $\alpha = 0.89$로 내용 타당도 및 내적 일관성이
확인되었으나, 대규모 패널에 의한 일반화 가능성은 추가 검증이 필요하다.
특히 산업별, 규모별 조직의 다양성을 반영한 전문가 패널 확장이
후속 연구에서 수행되어야 한다.

\chg{단일 $\kappa$ 값의 적용 측면에서,}
본 연구는 \jrev{\chg{\citet{lee2026drto}}에서 \chg{도출한}} $\kappa = 1.2$의 단일 곡률 매개변수를 전체 조건에 적용하였다.
\jrev{본 연구는 \frev{분포 환경 적합성} 종합 명제(\secref{sec:ch4_hypothesis_summary})를 통해 9개 가설(H2--H10) PASS의 통합 해석으로 $\kappa = 1.2$의 C2/C3 조건 \frev{운영 적합성을 시사하였으나},}
$\kappa$의 조건별 또는 산업별 차별화\jrev{, 그리고 환경 무관 다기준 종합점수의 재정의를 통한 직접 격자 탐색}이 모델의 적응력을 향상시킬 수 있으며,
이는 후속 연구의 과제이다.

\fdel{\noindent\textbf{시간적 역동성의 미반영.}
현재 $\DRTO$ 모델은 특정 시점의 관측성과 불확실성을 입력으로 하는
정적(static) 산출 구조이다. 실제 복구 과정에서 관측성과 불확실성은
시간에 따라 변화하므로, 시간적 역동성을 반영하는
동적(dynamic) 갱신 모델로의 확장이 필요하다.}
